\newglossaryentry{gls-AI}{
  type=fachwort,
  name=AI,
  first={Artificial Intelligence (AI)},
  text={AI},
  description={Artificial Intelligence - Englische Bezeichnung für KI}
}

\newglossaryentry{gls-ALSA}{
  type=fachwort,
  name=ALSA,
  first={Advanced Linux Sound Architecture (ALSA)},
  text={ALSA},
  description={Advanced Linux Sound Architecture, ein Framework im Linux-Kernel zur Verwaltung und Steuerung von Audio-Hardware}
}

\newglossaryentry{gls-API}{
  type=fachwort,
  name=API,
  first={Application Programming Interface (API)},
  text={API},
  description={Application Programming Interface, eine Schnittstelle für die Kommunikation zwischen Softwarekomponenten}
}

\newglossaryentry{gls-ARI}{
  type=fachwort,
  name=ARI,
  first={Asterisk REST Interface (ARI)},
  text={ARI},
  description={Asterisk REST Interface, eine REST-basierte Programmierschnittstelle zur Steuerung von Asterisk-Telefonsystemen und zur Implementierung eigener Kommunikationslogik}
}

\newglossaryentry{gls-ASR}{
  type=fachwort,
  name=ASR,
  first={Automatic Speech Recognition (ASR)},
  text={ASR},
  description={Automatic Speech Recognition, die automatische Erkennung gesprochener Sprache durch Maschinen}
}

\newglossaryentry{gls-AVV}{
  type=fachwort,
  name=AVV,
  first={Auftragsverarbeitungsvertrag (AVV)},
  text={AVV},
  description={Auftragsverarbeitungsvertrag, ein nach Art.28 DSGVO vorgeschriebener Vertrag zwischen Verantwortlichem und Auftragsverarbeiter, der Zweck, Art und Umfang der Datenverarbeitung sowie technische und organisatorische Ma\ss nahmen zum Schutz personenbezogener Daten regelt}
}

\newglossaryentry{gls-BPO}{
  type=fachwort,
  name=BPO,
  first={Business Process Outsourcing (BPO)},
  text={BPO},
  description={Auslagerung von Geschäftsprozessen an externe Dienstleister, um Kosten zu senken, Effizienz zu steigern oder sich auf Kernkompetenzen zu konzentrieren}
}

\newglossaryentry{gls-CDR}{
  type=fachwort,
  name=CDR,
  first={Call Detail Record (CDR)},
  text={CDR},
  description={Call Detail Record, ein Datensatz mit Informationen über einen durchgeführten Telefonanruf (Dauer, Teilnehmer, Uhrzeit, etc.)}
}

\newglossaryentry{gls-CLI}{
  type=fachwort,
  name=CLI,
  first={Command-Line Interface (CLI)},
  text={CLI},
  description={Command-Line Interface, eine textbasierte Benutzeroberfläche zur Bedienung von Programmen}
}

\newglossaryentry{gls-CPU}{
  type=fachwort,
  name=CPU,
  first={Central Processing Unit (CPU)},
  text={CPU},
  description={Central Processing Unit, der Hauptprozessor eines Computers, der Befehle ausführt}
}

\newglossaryentry{gls-DACH}{
  type=fachwort,
  name=DACH,
  first={Deutschland, Österreich und Schweiz (DACH)},
  text={DACH},
  description={Abkürzung für Deutschland, Österreich und die Schweiz im wirtschaftlichen und geografischen Kontext}
}

\newglossaryentry{gls-DID}{
  type=fachwort,
  name=DID,
  first={Direct Inward Dialing (DID)},
  text={DID},
  description={Direct Inward Dialing, ein Telekommunikationsverfahren, bei dem externe Anrufe direkt an eine interne Nebenstelle einer Telefonanlage weitergeleitet werden, ohne dass eine Vermittlung über eine zentrale Rufannahme erforderlich ist}
}

\newglossaryentry{gls-DPA}{
  type=fachwort,
  name=DPA,
  first={Data Processing Agreement (DPA)},
  text={DPA},
  description={Data Processing Agreement, ein vertragliches Dokument, das die Bedingungen der Datenverarbeitung zwischen Auftraggeber und Auftragsverarbeiter regelt}
}

\newglossaryentry{gls-DPIA}{
  type=fachwort,
  name=DPIA,
  first={Data Protection Impact Assessment (DPIA)},
  text={DPIA},
  description={Datenschutz-Folgenabschätzung anhand Art.~35 DSGVO zur systematischen Bewertung der Risiken einer Datenverarbeitung für die Rechte und Freiheiten betroffener Personen}
}

\newglossaryentry{gls-DSGVO}{
  type=fachwort,
  name=DSGVO,
  first={Datenschutz-Grundverordnung (DSGVO)},
  text={DSGVO},
  description={Datenschutz-Grundverordnung, eine EU-Verordnung zur Regelung des Schutzes personenbezogener Daten und zur Gewährleistung der informationellen Selbstbestimmung}
}

\newglossaryentry{gls-DTLS}{
  type=fachwort,
  name=DTLS,
  first={Datagram Transport Layer Security (DTLS)},
  text={DTLS},
  description={Datagram Transport Layer Security, ein Sicherheitsprotokoll für UDP-basierte Kommunikation}
}

\newglossaryentry{gls-DTMF}{
  type=fachwort,
  name=DTMF,
  first={Dual-Tone Multi-Frequency (DTMF)},
  text={DTMF},
  description={Dual-Tone Multi-Frequency, ein Verfahren zur Signalübertragung von Tastendrücken bei Telefonen, bei dem jede Taste durch eine Kombination aus zwei Frequenzen codiert wird}
}

\newglossaryentry{gls-E2E}{
  type=fachwort,
  name=E2E,
  first={End-to-End (E2E)},
  text={E2E},
  description={End-to-End, eine Bezeichnung für die Gesamtlatenz von Anfang bis Ende eines Prozesses oder einer Kommunikation}
}

\newglossaryentry{gls-EIA}{
  type=fachwort,
  name=EIA,
  first={Ethical Impact Assessment (EIA)},
  text={EIA},
  description={Systematische Bewertung der ethischen Auswirkungen von Technologien oder Systemen, insbesondere im Hinblick auf gesellschaftliche, rechtliche und moralische Konsequenzen}
}

\newglossaryentry{gls-ETSI}{
  type=fachwort,
  name=ETSI,
  first={European Telecommunications Standards Institute (ETSI)},
  text={ETSI},
  description={European Telecommunications Standards Institute, ein Standardisierungsgremium für Telekommunikation und Netzwerke}
}

\newglossaryentry{gls-GDPR}{
  type=fachwort,
  name=GDPR,
  first={General Data Protection Regulation (GDPR)},
  text={GDPR},
  description={Englische Bezeichnung der Datenschutz-Grundverordnung (DSGVO), die den rechtlichen Rahmen für die Verarbeitung personenbezogener Daten innerhalb der Europäischen Union festlegt}
}

\newglossaryentry{gls-HTTP}{
  type=fachwort,
  name=HTTP,
  first={Hypertext Transfer Protocol (HTTP)},
  text={HTTP},
  description={Hypertext Transfer Protocol, ein zustandsloses Anwendungsprotokoll zur Übertragung von Daten im World Wide Web}
}

\newglossaryentry{gls-HTTPS}{
  type=fachwort,
  name=HTTPS,
  first={Hypertext Transfer Protocol Secure (HTTPS)},
  text={HTTPS},
  description={Hypertext Transfer Protocol Secure, ein Protokoll zur sicheren Übertragung von Daten im Internet, das Verschlüsselung (SSL/TLS) und Authentifizierung von Webseiten bietet}
}

\newglossaryentry{gls-ICE}{
  type=fachwort,
  name=ICE,
  first={Interactive Connectivity Establishment (ICE)},
  text={ICE},
  description={Interactive Connectivity Establishment, ein Protokoll zur Überwindung von Firewall- und NAT-Hindernissen}
}

\newglossaryentry{gls-IETF}{
  type=fachwort,
  name=IETF,
  first={Internet Engineering Task Force (IETF)},
  text={IETF},
  description={Internet Engineering Task Force, eine internationale Organisation, die offene Standards für das Internet entwickelt, insbesondere Protokolle und Architekturen}
}

\newglossaryentry{gls-ISO}{
  type=fachwort,
  name=ISO,
  first={International Organization for Standardization (ISO)},
  text={ISO},
  description={International Organization for Standardization, eine internationale Organisation zur Entwicklung und Veröffentlichung von Normen für verschiedene Technologien, Prozesse und Qualitätsstandards}
}

\newglossaryentry{gls-ITU-T}{
  type=fachwort,
  name=ITU-T,
  first={International Telecommunication Union (ITU-T)},
  text={ITU-T},
  description={International Telecommunication Union - Standardisierungsgremium für Telekommunikation}
}

\newglossaryentry{gls-IVR}{
  type=fachwort,
  name=IVR,
  first={Interactive Voice Response (IVR)},
  text={IVR},
  description={Interactive Voice Response, ein System für automatisierte Sprachmenüs und Sprachverarbeitung}
}

\newglossaryentry{gls-JSON}{
  type=fachwort,
  name=JSON,
  first={JavaScript Object Notation (JSON)},
  text={JSON},
  description={JavaScript Object Notation, ein text-basiertes Datenformat für den Austausch strukturierter Daten}
}

\newglossaryentry{gls-JWT}{
  type=fachwort,
  name=JWT,
  first={JSON Web Token (JWT)},
  text={JWT},
  description={JSON Web Token, ein offener Standard (RFC 7519) zur sicheren Übertragung von Claims zwischen Parteien, häufig verwendet für Authentifizierung und Autorisierung in Webanwendungen}
}

\newglossaryentry{gls-KI}{
  type=fachwort,
  name=KI,
  first={Künstliche Intelligenz (KI)},
  text={KI},
  description={Künstliche Intelligenz bezeichnet die Fähigkeit von Maschinen, menschenähnliche kognitive Leistungen wie Lernen, Problemlösen und Sprachverstehen auszuführen}
}

\newglossaryentry{gls-KMS}{
  type=fachwort,
  name=KMS,
  first={Key Management Service (KMS)},
  text={KMS},
  description={Key Management Service, ein Dienst zur sicheren Verwaltung, Speicherung und Rotation kryptografischer Schlüssel, der häufig in Cloud- und Container-Umgebungen eingesetzt wird}
}

\newglossaryentry{gls-KMU}{
  type=fachwort,
  name=KMU,
  first={kleine und mittlere Unternehmen (KMU)},
  text={KMU},
  description={kleine und mittelständische Unternehmen}
}

\newglossaryentry{gls-LAN}{
  type=fachwort,
  name=LAN,
  first={Local Area Network (LAN)},
  text={LAN},
  description={Local Area Network, ein räumlich begrenztes Rechnernetzwerk, das Geräte innerhalb eines lokalen Bereichs wie eines Gebäudes oder Campus verbindet}
}

\newglossaryentry{gls-LLM}{
  type=fachwort,
  name=LLM,
  first={Large Language Model (LLM)},
  text={LLM},
  description={Large Language Model, ein auf gro\ss en Textkorpora trainiertes KI-Sprachmodell, das natürliche Sprache generieren oder verstehen kann}
}

\newglossaryentry{gls-MANO}{
  type=fachwort,
  name=MANO,
  first={Management and Orchestration (MANO)},
  text={MANO},
  description={Eine Architektur- und Funktionsschicht im NFV-Umfeld zur Verwaltung, Orchestrierung und Überwachung von virtualisierten Netzwerkfunktionen (VNFs) sowie der zugrunde liegenden Infrastruktur}
}

\newglossaryentry{gls-MEC}{
  type=fachwort,
  name=MEC,
  first={Multi-access Edge Computing (MEC)},
  text={MEC},
  description={Multi-access Edge Computing, ein Architekturkonzept, bei dem Rechen- und Speicherkapazitäten nahe am Netzwerkrand bereitgestellt werden, um Latenzen zu reduzieren und zeitkritische Anwendungen zu unterstützen}
}

\newglossaryentry{gls-MOS}{
  type=fachwort,
  name=MOS,
  first={Mean Opinion Score (MOS)},
  text={MOS},
  description={Mean Opinion Score, eine Metrik zur Bewertung der Audioqualität auf einer Skala}
}

\newglossaryentry{gls-NAT}{
  type=fachwort,
  name=NAT,
  first={Network Address Translation (NAT)},
  text={NAT},
  description={Network Address Translation, ein Verfahren zur Umwandlung von IP-Adressen in Netzwerkpaketen}
}

\newglossaryentry{gls-NDJSON}{
  type=fachwort,
  name=NDJSON,
  first={Newline Delimited JSON (NDJSON)},
  text={NDJSON},
  description={Newline Delimited JSON, ein Datenformat zur sequentiellen Darstellung von JSON-Objekten, bei dem jedes Objekt in einer eigenen Zeile steht, wodurch sich das Format besonders für Streaming, Log-Dateien und die Verarbeitung gro\ss er Datenmengen eignet}
}

\newglossaryentry{gls-NFV}{
  type=fachwort,
  name=NFV,
  first={Network Function Virtualization (NFV)},
  text={NFV},
  description={Network Function Virtualization, ein Paradigma zur Ablösung dedizierter Hardware durch softwarebasierte, containerisierte Netzwerkdienste}
}

\newglossaryentry{gls-NFVI}{
  type=fachwort,
  name=NFVI,
  first={Network Function Virtualization Infrastructure (NFVI)},
  text={NFVI},
  description={Network Function Virtualization Infrastructure, die physische Hardware- und Software-Ressourcenschicht, auf der VNFs laufen}
}

\newglossaryentry{gls-NFVO}{
  type=fachwort,
  name=NFVO,
  first={Network Function Virtualization Orchestration (NFVO)},
  text={NFVO},
  description={Network Function Virtualization Orchestration, die Verwaltungs- und Orchestrierungsschicht für VNFs und deren Lifecycle-Management}
}

\newglossaryentry{gls-NLP}{
  type=fachwort,
  name=NLP,
  first={Natural Language Processing (NLP)},
  text={NLP},
  description={Natural Language Processing bezeichnet die automatische Verarbeitung, Analyse und Generierung menschlicher Sprache durch Computer, z.B. für Textverständnis, Sprachmodellierung oder Dialogsysteme}
}

\newglossaryentry{gls-OSS}{
  type=fachwort,
  name=OSS,
  first={Open Source Software (OSS)},
  text={OSS},
  description={Open Source Software, Software mit öffentlich verfügbarem Quellcode, die frei verändert und verteilt werden darf}
}

\newglossaryentry{gls-PABX}{
  type=fachwort,
  name=PABX,
  first={Private Automatic Branch Exchange (PABX)},
  text={PABX},
  description={Private Automatic Branch Exchange, eine private, automatische Nebenstellenanlage zur Vermittlung interner und externer Sprachverbindungen in Unternehmensnetzwerken. Auch als klassische Telefonanlage oder Private Branch Exchange (PBX) bekannt}
}

\newglossaryentry{gls-PBX}{
  type=fachwort,
  name=PBX,
  first={Private Branch Exchange (PBX)},
  text={PBX},
  description={Private Branch Exchange, eine Telefonanlage für interne und externe Sprachkommunikation}
}

\newglossaryentry{gls-PCM}{
  type=fachwort,
  name=PCM,
  first={Pulse Code Modulation (PCM)},
  text={PCM},
  description={Pulse Code Modulation, ein Verfahren zur digitalen Darstellung analoger Signale durch Abtastung, Quantisierung und Codierung, häufig verwendet in der digitalen Sprach- und Audioübertragung}
}

\newglossaryentry{gls-PCM16}{
  type=fachwort,
  name=PCM16,
  first={Pulse Code Modulation 16 Bit (PCM16)},
  text={PCM16},
  description={Pulse Code Modulation mit 16 Bit Auflösung, ein unkomprimiertes Audioformat zur digitalen Darstellung von Audiosignalen, das häufig in der Sprachverarbeitung und Echtzeitkommunikation verwendet wird}
}

\newglossaryentry{gls-PESQ}{
  type=fachwort,
  name=PESQ,
  first={Perceptual Evaluation of Speech Quality (PESQ)},
  text={PESQ},
  description={Perceptual Evaluation of Speech Quality, ein objektives Verfahren zur Bewertung der Sprachqualität durch Vergleich eines Referenzsignals mit einem gestörten Signal, standardisiert in ITU-T P.862 und weit verbreitet zur Bewertung von Sprachübertragungssystemen}
}

\newglossaryentry{gls-PII}{
  type=fachwort,
  name=PII,
  first={Personally Identifiable Information (PII)},
  text={PII},
  description={Personally Identifiable Information, persönliche Daten, die zur Identifikation einer Person verwendet werden können}
}

\newglossaryentry{gls-PJSIP}{
  type=fachwort,
  name=PJSIP,
  first={PJSIP (PJSIP)},
  text={PJSIP},
  description={Open-Source-Bibliothek für SIP-basierte Sprach- und Videoübertragung, die Funktionen für VoIP, Instant Messaging und NAT-Traversierung bereitstellt}
}

\newglossaryentry{gls-POLQA}{
  type=fachwort,
  name=POLQA,
  first={Perceptual Objective Listening Quality Assessment (POLQA)},
  text={POLQA},
  description={Perceptual Objective Listening Quality Assessment, ein objektives Messverfahren zur Bewertung der wahrgenommenen Sprachqualität in Telekommunikationssystemen, standardisiert in ITU-T P.863 und ausgelegt für moderne Sprachdienste wie VoIP, HD-Voice und Mobilfunknetze}
}

\newglossaryentry{gls-PSTN}{
  type=fachwort,
  name=PSTN,
  first={Public Switched Telephone Network (PSTN)},
  text={PSTN},
  description={Public Switched Telephone Network, das klassische öffentliche Telefonnetz mit Sprachübertragung}
}

\newglossaryentry{gls-RAG}{
  type=fachwort,
  name=RAG,
  first={Retrieval-Augmented Generation (RAG)},
  text={RAG},
  description={Retrieval-Augmented Generation, ein Ansatz in der natürlichen Sprachverarbeitung, bei dem ein generatives Sprachmodell mit einer externen Wissensquelle kombiniert wird, um Antworten auf Basis abgerufener, kontextrelevanter Informationen zu erzeugen}
}

\newglossaryentry{gls-RAM}{
  type=fachwort,
  name=RAM,
  first={Random Access Memory (RAM)},
  text={RAM},
  description={Random Access Memory, der Arbeitsspeicher eines Computers für temporäre Datenablage}
}

\newglossaryentry{gls-RBAC}{
  type=fachwort,
  name=RBAC,
  first={Role-Based Access Control (RBAC)},
  text={RBAC},
  description={Role-Based Access Control, ein Zugriffskontrollmodell, das Berechtigungen auf Basis von Benutzerrollen vergibt}
}

\newglossaryentry{gls-REST}{
  type=fachwort,
  name=REST,
  first={Representational State Transfer (REST)},
  text={REST},
  description={Representational State Transfer, ein Architekturstil für die Gestaltung von Web-APIs}
}

\newglossaryentry{gls-RFC}{
  type=fachwort,
  name=RFC,
  first={Request for Comments (RFC)},
  text={RFC},
  description={Request for Comments, eine von der Internet Engineering Task Force veröffentlichte Dokumentenreihe zur Spezifikation, Beschreibung und Diskussion von Internetstandards und -protokollen}
}

\newglossaryentry{gls-RMS}{
  type=fachwort,
  name=RMS,
  first={Root Mean Square (RMS)},
  text={RMS},
  description={Root Mean Square, ein mathematisches Ma\ss zur Bestimmung des quadratischen Mittelwerts eines Signals, häufig verwendet zur Beschreibung der effektiven Stärke von Wechselspannungen oder Strömen}
}

\newglossaryentry{gls-RTC}{
  type=fachwort,
  name=RTC,
  first={Real-Time Communication (RTC)},
  text={RTC},
  description={Real-Time Communication, Sammelbegriff für Technologien zur zeitkritischen Übertragung von Audio-, Video- und Datenströmen in Echtzeit über Netzwerke}
}

\newglossaryentry{gls-RTCP}{
  type=fachwort,
  name=RTCP,
  first={Real-Time Control Protocol (RTCP)},
  text={RTCP},
  description={Ein Begleitprotokoll zu RTP, das für die Überwachung und Kontrolle von Echtzeit-Medienströmen zuständig ist}
}

\newglossaryentry{gls-RTP}{
  type=fachwort,
  name=RTP,
  first={Real-time Transport Protocol (RTP)},
  text={RTP},
  description={Real-time Transport Protocol, das Protokoll zur Übertragung von Audio und Videodaten in Echtzeit}
}

\newglossaryentry{gls-SCC}{
  type=fachwort,
  name=SCC,
  first={Standard Contractual Clauses (SCC)},
  text={SCC},
  description={Standard Contractual Clauses, von der EU bereitgestellte Standardvertragsmuster für die sichere Datenübermittlung in Drittstaaten}
}

\newglossaryentry{gls-SDK}{
  type=fachwort,
  name=SDK,
  first={Software Development Kit (SDK)},
  text=SDK,
  description={Software Development Kit, eine Sammlung von Werkzeugen und Bibliotheken zur Entwicklung von Anwendungen}
}

\newglossaryentry{gls-SFU}{
  type=fachwort,
  name=SFU,
  first={Selective Forwarding Unit (SFU)},
  text=SFU,
  description={Selective Forwarding Unit, ein Media-Server für WebRTC, der Audio- und Videodatenströme zwischen Teilnehmern weiterleitet}
}

\newglossaryentry{gls-SIP}{
  type=fachwort,
  name=SIP,
  first={Session Initiation Protocol (SIP)},
  text={SIP},
  description={Session Initiation Protocol, ein Standardprotokoll zur Signalisierung und Verwaltung von VoIP-Anrufen}
}

\newglossaryentry{gls-SLA}{
  type=fachwort,
  name=SLA,
  first={Service Level Agreement (SLA)},
  text=SLA,
  description={Service Level Agreement, eine vertragliche Vereinbarung, die die erwartete Servicequalität (z.B. Verfügbarkeit, Latenz) definiert}
}

\newglossaryentry{gls-SPA}{
  type=fachwort,
  name=SPA,
  first={Single-Page Application (SPA)},
  text=SPA,
  description={Single-Page Application, eine Webanwendung, die in einem Single Document geladen wird und dynamisch Inhalte aktualisiert}
}

\newglossaryentry{gls-SRTP}{
  type=fachwort,
  name=SRTP,
  first={Secure Real-time Transport Protocol (SRTP)},
  text=SRTP,
  description={Secure Real-time Transport Protocol, eine sichere Variante von RTP mit Verschlüsselung}
}

\newglossaryentry{gls-STT}{
  type=fachwort,
  name=STT,
  first={Speech-to-Text (STT)},
  text=STT,
  description={Speech-to-Text, die Umwandlung von Sprache in geschriebenen Text}
}

\newglossaryentry{gls-STUN}{
  type=fachwort,
  name=STUN,
  first={Session Traversal Utilities for NAT (STUN)},
  text=STUN,
  description={Session Traversal Utilities for NAT, ein Protokoll zur Ermittlung öffentlicher IP-Adressen und Ports zur Unterstützung der Kommunikation durch NAT-Gateways}
}

\newglossaryentry{gls-TCP}{
  type=fachwort,
  name=TCP,
  first={Transmission Control Protocol (TCP)},
  text=TCP,
  description={Transmission Control Protocol, ein verbindungsorientiertes Transportprotokoll der Internetprotokollfamilie, das eine zuverlässige, geordnete und fehlerfreie Datenübertragung gewährleistet}
}

\newglossaryentry{gls-TLS}{
  type=fachwort,
  name=TLS,
  first={Transport Layer Security (TLS)},
  text=TLS,
  description={Transport Layer Security, ein Kryptografieprotokoll für sichere Kommunikation über Netzwerke}
}

\newglossaryentry{gls-TTL}{
  type=fachwort,
  name=TTL,
  first={Time-To-Live (TTL)},
  text=TTL,
  description={Time-To-Live, ein Parameter in Netzwerkprotokollen (insbesondere IP), der die maximale Anzahl von Hops (Sprüngen) begrenzt, die ein Paket durchlaufen kann, bevor es verworfen wird}
}

\newglossaryentry{gls-TTS}{
  type=fachwort,
  name=TTS,
  first={Text-to-Speech (TTS)},
  text=TTS,
  description={Text-to-Speech, die Umwandlung von geschriebenem Text in gesprochene Sprache}
}

\newglossaryentry{gls-TURN}{
  type=fachwort,
  name=TURN,
  first={Traversal Using Relays around NAT (TURN)},
  text=TURN,
  description={Traversal Using Relays around NAT, ein Protokoll zur relay-basierten Weiterleitung von Medienströmen, wenn direkte Verbindungen durch NAT oder Firewalls nicht möglich sind}
}

\newglossaryentry{gls-UDP}{
  type=fachwort,
  name=UDP,
  first={User Datagram Protocol (UDP)},
  text=UDP,
  description={User Datagram Protocol, ein verbindungsloses Transportprotokoll der Internetprotokollfamilie, das eine schnelle, aber nicht garantierte Übertragung von Daten ohne Reihenfolge oder Zustellsicherung ermöglicht}
}

\newglossaryentry{gls-UI}{
  type=fachwort,
  name=UI,
  first={User Interface (UI)},
  text=UI,
  description={User Interface, die Benutzeroberfläche einer Anwendung, über die Benutzer mit dem System interagieren}
}

\newglossaryentry{gls-VAD}{
  type=fachwort,
  name=VAD,
  first={Voice Activity Detection (VAD)},
  text=VAD,
  description={Voice Activity Detection, die Erkennung von Sprachaktivität in Audiodaten}
}

\newglossaryentry{gls-VLAN}{
  type=fachwort,
  name=VLAN,
  first={Virtual Local Area Network (VLAN)},
  text=VLAN,
  description={Virtual Local Area Network, ein logisches Teilnetz zur Segmentierung von Netzwerken auf Layer 2, das zur Trennung von Datenverkehr und zur Erhöhung von Sicherheit und Übersichtlichkeit eingesetzt wird}
}

\newglossaryentry{gls-VM}{
  type=fachwort,
  name=VM,
  first={Virtuelle Maschine (VM)},
  text=VM,
  description={Virtuelle Maschine, eine softwarebasierte Emulation eines Computersystems, die es ermöglicht, Betriebssysteme und Anwendungen isoliert auf derselben Hardware auszuführen}
}

\newglossaryentry{gls-VNF}{
  type=fachwort,
  name=VNF,
  first={Virtualized Network Function (VNF)},
  text=VNF,
  description={Virtualized Network Function, eine softwarebasierte Implementierung einer Netzwerkfunktion, die auf Standard-Hardware läuft}
}

\newglossaryentry{gls-YAML}{
  type=fachwort,
  name=YAML,
  first={YAML Ain't Markup Language (YAML)},
  text=YAML,
  description={YAML Ain’t Markup Language, ein menschenlesbares Datenserialisierungsformat zur strukturierten Darstellung von Konfigurationsdaten, das durch Einrückungen anstelle von Klammern oder Tags arbeitet}
}

\newglossaryentry{gls-Asterisk}{
  type=fachwort,
  name=Asterisk,
  description={Open-Source-Software-PBX, die Signalisierung, Rufaufbau und Medienverarbeitung für VoIP- und klassische Telefonsysteme bereitstellt}
}
\glsadd{gls-Asterisk}

\newglossaryentry{gls-AudioDec}{
  type=fachwort,
  name=AudioDec,
  description={Neuraler Audio-Codec-Ansatz mit sehr niedriger Latenz, der für Echtzeit-Anwendungen optimiert ist}
}
\glsadd{gls-AudioDec}

\newglossaryentry{gls-Aufzeichnungs-Opt-in}{
  type=fachwort,
  name=Aufzeichnungs-Opt-in,
  description={Die ausdrückliche Zustimmung eines Benutzers oder Gesprächspartners zur Aufzeichnung von Audio-, Video- oder Transaktionsdaten, die für Compliance-, Qualitäts- oder Analysezwecke erhoben werden}
}
\glsadd{gls-Aufzeichnungs-Opt-in}

\newglossaryentry{gls-Bridge}{
  type=fachwort,
  name=Bridge,
  description={Netzwerk- oder Softwarekomponente, die zwei unterschiedliche Netzwerke, Protokolle oder Systeme verbindet und Datenaustausch ermöglicht, z.B. SIP-WebRTC-Bridge}
}
\glsadd{gls-Bridge}

\newglossaryentry{gls-Codec}{
  type=fachwort,
  name=Codec,
  description={Encoder-Decoder-Verfahren zur Digitalisierung, Kompression und Dekomprimierung von Audio- oder Videodaten. Beispiele: G.711 (klassische Telefonie), Opus (WebRTC), PCM16 (Spracherkennung)}
}
\glsadd{gls-Codec}

\newglossaryentry{gls-Container}{
  type=fachwort,
  name=Container,
  description={Eine leichtgewichtige virtualisierte Laufzeitumgebung, in der Software und ihre Abhängigkeiten isoliert ausgeführt werden}
}
\glsadd{gls-Container}

\newglossaryentry{gls-Whisper}{
  type=fachwort,
  name=Whisper,
  description={OpenAI-Spracherkennungsmodell für mehrsprachige Speech-to-Text-Verarbeitung, verfügbar als Cloud-API und als lokal ausführbares Modell}
}
\glsadd{gls-Whisper}

\newglossaryentry{gls-Whitelist}{
  type=fachwort,
  name=Whitelist,
  description={Liste von erlaubten IP-Adressen, Benutzern oder anderen Entities. Nur Einträge auf der Whitelist werden akzeptiert, alle anderen werden blockiert}
}
\glsadd{gls-Whitelist}

\newglossaryentry{gls-Baresip}{
  type=fachwort,
  name=Baresip,
  description={Open-Source SIP User Agent für experimentelle Tests und grundlegende SIP-Registrierung, der für Protokollkompatibilitätsprüfungen eingesetzt wird}
}
\glsadd{gls-Baresip}

\newglossaryentry{gls-LiveKit}{
  type=fachwort,
  name=LiveKit,
  description={Open-Source WebRTC-Framework und Media-Server für Echtzeitkommunikation, das Raumerstellung, Audiostream-Verarbeitung und KI-Integration unterstützt}
}
\glsadd{gls-LiveKit}

\newglossaryentry{gls-Media-Proxy}{
  type=fachwort,
  name=Media-Proxy,
  description={Netzwerkkomponente zur Weiterleitung, Anpassung und Transkodierung von Audio- und Videoströmen zwischen unterschiedlichen Protokollen und Codecs}
}
\glsadd{gls-Media-Proxy}

\newglossaryentry{gls-RTP-Bridge}{
  type=fachwort,
  name=RTP-Bridge,
  description={Komponente zur Weiterleitung, Anpassung und Transkodierung von RTP-Medienströmen zwischen verschiedenen Diensten und Protokollen}
}
\glsadd{gls-RTP-Bridge}

\newglossaryentry{gls-Hypervisor}{
  type=fachwort,
  name=Hypervisor,
  description={Virtualisierungssoftware, die es ermöglicht, mehrere virtuelle Maschinen auf einem physischen Host-System auszuführen und deren Ressourcenzugriff zu verwalten}
}
\glsadd{gls-Hypervisor}

\newglossaryentry{gls-Docker}{
  type=fachwort,
  name=Docker,
  description={Containerisierungsplattform zur Erstellung, Verwaltung und Ausführung von Anwendungen in isolierten Container-Umgebungen}
}
\glsadd{gls-Docker}

\newglossaryentry{gls-Docker-Image}{
  type=fachwort,
  name=Docker-Image,
  description={Schichtenbasierte, unveränderliche Vorlage zur Erstellung von Docker-Containern, die alle benötigten Abhängigkeiten, Bibliotheken und Konfigurationen enthält}
}
\glsadd{gls-Docker-Image}

\newglossaryentry{gls-Kubernetes}{
  type=fachwort,
  name=Kubernetes,
  description={Open-Source-Orchestrierungsplattform für die Automatisierung von Bereitstellung, Skalierung und Verwaltung containerisierter Anwendungen}
}
\glsadd{gls-Kubernetes}

\newglossaryentry{gls-Kamailio}{
  type=fachwort,
  name=Kamailio,
  description={Open-Source SIP-Server und Proxy für hochperformante VoIP-Telefonie, der für Lastverteilung und SIP-Routing eingesetzt wird}
}
\glsadd{gls-Kamailio}

\newglossaryentry{gls-Dialplan}{
  type=fachwort,
  name=Dialplan,
  description={Regelwerk in einer PBX-Software zur Definition von Anrufverarbeitungslogik, Routing-Entscheidungen und Aktionsabläufen}
}
\glsadd{gls-Dialplan}

\newglossaryentry{gls-Jitter-Puffer}{
  type=fachwort,
  name=Jitter-Puffer,
  description={Zwischenspeicher zur Kompensation von Verzögerungsschwankungen (Jitter) bei der Übertragung von Medienpaketen in IP-Netzwerken, um eine gleichmäßige Wiedergabe zu gewährleisten}
}
\glsadd{gls-Jitter-Puffer}

\newglossaryentry{gls-Streaming}{
  type=fachwort,
  name=Streaming,
  description={Übertragungstechnik, bei der Daten kontinuierlich in Echtzeit gesendet und verarbeitet werden, ohne das vollständige Ergebnis abzuwarten}
}
\glsadd{gls-Streaming}

\newglossaryentry{gls-Token-LLM}{
  type=fachwort,
  name=Token,
  description={Kleinste Verarbeitungseinheit in Sprachmodellen, die Wörter, Wortteile oder Zeichen repräsentiert und bei der Textgenerierung verwendet wird}
}
\glsadd{gls-Token-LLM}

\newglossaryentry{gls-ConfigMaps}{
  type=fachwort,
  name=ConfigMaps,
  description={Kubernetes-Ressource zur Verwaltung von Konfigurationsdaten als Key-Value-Paare, die von Containern geladen werden können, ohne dass diese in Container-Images eingebettet werden müssen}
}
\glsadd{gls-ConfigMaps}

\newglossaryentry{gls-Secrets}{
  type=fachwort,
  name=Secrets,
  description={Kubernetes-Objekt zur sicheren Speicherung vertraulicher Daten wie Passwörter, API-Schlüssel und Zertifikate}
}
\glsadd{gls-Secrets}

\newglossaryentry{gls-Persistent-Volume}{
  type=fachwort,
  name=Persistent Volume,
  description={Kubernetes-Ressource zur Bereitstellung dauerhaften Speicherplatzes für Container, der über den Lebenszyklus einzelner Pods hinaus erhalten bleibt}
}
\glsadd{gls-Persistent-Volume}

\newglossaryentry{gls-ReplicaSet}{
  type=fachwort,
  name=ReplicaSet,
  description={Kubernetes-Controller zur Sicherstellung, dass eine definierte Anzahl identischer Pod-Replikas jederzeit ausgeführt wird}
}
\glsadd{gls-ReplicaSet}

\newglossaryentry{gls-Pod}{
  type=fachwort,
  name=Pod,
  description={Kleinste ausführbare Einheit in Kubernetes, die einen oder mehrere Container enthält und gemeinsam bereitgestellt wird}
}
\glsadd{gls-Pod}

\newglossaryentry{gls-Namespace}{
  type=fachwort,
  name=Namespace,
  description={Logische Trennung von Ressourcen innerhalb eines Kubernetes-Clusters zur Isolation verschiedener Anwendungen oder Teams}
}
\glsadd{gls-Namespace}

\newglossaryentry{gls-NetworkPolicy}{
  type=fachwort,
  name=NetworkPolicy,
  description={Kubernetes-Ressource zur Definition netzwerkbasierter Zugriffskontrollregeln zwischen Pods}
}
\glsadd{gls-NetworkPolicy}

\newglossaryentry{gls-Ingress-Controller}{
  type=fachwort,
  name=Ingress-Controller,
  description={Kubernetes-Komponente zur Verwaltung des externen HTTP/HTTPS-Zugriffs auf Services innerhalb eines Clusters}
}
\glsadd{gls-Ingress-Controller}

\newglossaryentry{gls-Least-Privilege-Prinzip}{
  type=fachwort,
  name=Least-Privilege-Prinzip,
  description={Sicherheitsprinzip, bei dem Benutzern und Prozessen nur die minimal notwendigen Berechtigungen erteilt werden, um ihre Aufgaben zu erfüllen}
}
\glsadd{gls-Least-Privilege-Prinzip}

\newglossaryentry{gls-Digest-Authentifizierung}{
  type=fachwort,
  name=Digest-Authentifizierung,
  description={HTTP-Authentifizierungsverfahren, bei dem Passwörter gehasht übertragen werden, um Abhören zu erschweren}
}
\glsadd{gls-Digest-Authentifizierung}

\newglossaryentry{gls-Horizontal-Pod-Autoscaler}{
  type=fachwort,
  name=Horizontal Pod Autoscaler,
  description={Kubernetes-Mechanismus zur automatischen Anpassung der Anzahl von Pod-Replicas basierend auf CPU-, Speicher- oder benutzerdefinierten Metriken}
}
\glsadd{gls-Horizontal-Pod-Autoscaler}

\newglossaryentry{gls-SIP-Trunk}{
  type=fachwort,
  name=SIP-Trunk,
  description={Virtuelle Verbindungsleitung zwischen einer PBX und einem VoIP-Provider zur Übertragung mehrerer gleichzeitiger Sprachkanäle über IP-Netzwerke}
}
\glsadd{gls-SIP-Trunk}

\newglossaryentry{gls-Intent-Erkennung}{
  type=fachwort,
  name=Intent-Erkennung,
  description={Verfahren zur automatischen Identifikation der Absicht oder des Ziels einer Nutzeräußerung in einem Dialog, häufig eingesetzt in Chatbots und Sprachassistenten}
}
\glsadd{gls-Intent-Erkennung}

\newglossaryentry{gls-Dialog-Management}{
  type=fachwort,
  name=Dialog-Management,
  description={Steuerungslogik für die Verwaltung des Gesprächsflusses in einem KI-basierten Dialog, die Kontexte verfolgt und Antworten koordiniert}
}
\glsadd{gls-Dialog-Management}

\newglossaryentry{gls-Task-Orchestrierung}{
  type=fachwort,
  name=Task-Orchestrierung,
  description={Koordination und Steuerung mehrerer zusammenhängender Aufgaben oder Prozesse in einem System}
}
\glsadd{gls-Task-Orchestrierung}

\newglossaryentry{gls-Geocodierung}{
  type=fachwort,
  name=Geocodierung,
  description={Prozess zur Umwandlung von Adressinformationen in geografische Koordinaten (Längen- und Breitengrad)}
}
\glsadd{gls-Geocodierung}

\newglossaryentry{gls-Prometheus}{
  type=fachwort,
  name=Prometheus,
  description={Open-Source-Monitoring-System zur Erfassung und Speicherung von Zeitreihendaten (Metriken) aus IT-Systemen}
}
\glsadd{gls-Prometheus}

\newglossaryentry{gls-Grafana}{
  type=fachwort,
  name=Grafana,
  description={Open-Source-Plattform zur Visualisierung und Analyse von Metriken, Logs und Traces aus verschiedenen Datenquellen}
}
\glsadd{gls-Grafana}

\newglossaryentry{gls-PostgreSQL}{
  type=fachwort,
  name=PostgreSQL,
  description={Open-Source relationales Datenbankmanagementsystem mit umfangreicher SQL-Unterstützung}
}
\glsadd{gls-PostgreSQL}

\newglossaryentry{gls-MySQL}{
  type=fachwort,
  name=MySQL,
  description={Weit verbreitetes Open-Source relationales Datenbankmanagementsystem}
}
\glsadd{gls-MySQL}

\newglossaryentry{gls-Monolithisches-Modell}{
  type=fachwort,
  name=Monolithisches Modell,
  description={Softwarearchitektur, bei der alle Funktionen und Komponenten in einer einzigen, eng gekoppelten Anwendung integriert sind}
}
\glsadd{gls-Monolithisches-Modell}

\newglossaryentry{gls-Mikroservice}{
  type=fachwort,
  name=Mikroservice,
  description={Architekturmuster, bei dem eine Anwendung in kleine, unabhängig deploybare Services aufgeteilt wird, die über definierte Schnittstellen kommunizieren}
}
\glsadd{gls-Mikroservice}

\newglossaryentry{gls-Hybrid-Ansatz}{
  type=fachwort,
  name=Hybrid-Ansatz,
  description={Architekturmodell, das Elemente monolithischer und mikroservice-basierter Ansätze kombiniert}
}
\glsadd{gls-Hybrid-Ansatz}

\newglossaryentry{gls-Barge-in}{
  type=fachwort,
  name=Barge-in,
  description={Funktion in Sprachdialogsystemen, die es Nutzern ermöglicht, die Systemausgabe zu unterbrechen und sofort zu antworten}
}
\glsadd{gls-Barge-in}

\newglossaryentry{gls-Edge-Deployment}{
  type=fachwort,
  name=Edge-Deployment,
  description={Bereitstellungsmodell, bei dem Anwendungen nahe am Netzwerkrand (z.B. auf lokalen Servern, Edge-Knoten oder IoT-Geräten) statt in zentralen Rechenzentren betrieben werden}
}
\glsadd{gls-Edge-Deployment}

\newglossaryentry{gls-Vendor-Lock-in}{
  type=fachwort,
  name=Vendor-Lock-in,
  description={Abhängigkeit von einem spezifischen Anbieter, bei der ein Wechsel zu alternativen Lösungen technisch oder wirtschaftlich erschwert wird}
}
\glsadd{gls-Vendor-Lock-in}

\newglossaryentry{gls-On-Premise}{
  type=fachwort,
  name=On-Premise,
  description={Betriebsmodell, bei dem Software und Infrastruktur auf eigenen Servern im Unternehmen oder lokalen Rechenzentren statt in der Cloud betrieben werden}
}
\glsadd{gls-On-Premise}

\newglossaryentry{gls-GitHub-Copilot}{
  type=fachwort,
  name=GitHub Copilot,
  description={KI-gestütztes Programmierwerkzeug zur automatischen Code-Vervollständigung und Generierung von Code-Vorschlägen}
}
\glsadd{gls-GitHub-Copilot}

\newglossaryentry{gls-Next-js}{
  type=fachwort,
  name=Next.js,
  description={React-Framework für serverseitiges Rendering und statische Website-Generierung}
}
\glsadd{gls-Next-js}

\newglossaryentry{gls-React}{
  type=fachwort,
  name=React,
  description={JavaScript-Bibliothek zur Entwicklung von Benutzeroberflächen mit komponentenbasierter Architektur}
}
\glsadd{gls-React}

\newglossaryentry{gls-Node-js}{
  type=fachwort,
  name=Node.js,
  description={JavaScript-Laufzeitumgebung für serverseitige Anwendungen}
}
\glsadd{gls-Node-js}

\newglossaryentry{gls-npm}{
  type=fachwort,
  name=npm,
  description={Package Manager für JavaScript und Node.js zur Verwaltung von Abhängigkeiten}
}
\glsadd{gls-npm}

\newglossaryentry{gls-pip}{
  type=fachwort,
  name=pip,
  description={Package Installer für Python zur Installation und Verwaltung von Python-Paketen}
}
\glsadd{gls-pip}

\newglossaryentry{gls-Docker-Desktop}{
  type=fachwort,
  name=Docker Desktop,
  description={Desktop-Anwendung zur lokalen Entwicklung und Verwaltung von Docker-Containern unter Windows und macOS}
}
\glsadd{gls-Docker-Desktop}

\newglossaryentry{gls-Docker-Compose}{
  type=fachwort,
  name=Docker-Compose,
  description={Werkzeug zur Definition und Verwaltung von Multi-Container-Docker-Anwendungen mittels YAML-Konfigurationsdateien}
}
\glsadd{gls-Docker-Compose}

\newglossaryentry{gls-Bridge-Netzwerk}{
  type=fachwort,
  name=Bridge-Netzwerk,
  description={Docker-Netzwerkmodus, bei dem Container über ein virtuelles Netzwerk miteinander kommunizieren}
}
\glsadd{gls-Bridge-Netzwerk}

\newglossaryentry{gls-Redis}{
  type=fachwort,
  name=Redis,
  description={Open-Source In-Memory-Datenbank für schnelle Datenzugriffe, häufig für Caching und Session-Verwaltung eingesetzt}
}
\glsadd{gls-Redis}

\newglossaryentry{gls-Observer-Pattern}{
  type=fachwort,
  name=Observer-Pattern,
  description={Entwurfsmuster, bei dem Objekte automatisch über Zustandsänderungen anderer Objekte benachrichtigt werden}
}
\glsadd{gls-Observer-Pattern}

\newglossaryentry{gls-Hook}{
  type=fachwort,
  name=Hook,
  description={React-Funktion zur Nutzung von Zustandsverwaltung und Seiteneffekten in funktionalen Komponenten}
}
\glsadd{gls-Hook}

\newglossaryentry{gls-Server-Side-Token-Pattern}{
  type=fachwort,
  name=Server-Side Token Pattern,
  description={Architekturmuster, bei dem Authentifizierungs-Tokens serverseitig generiert werden, um Geheimnisse vor dem Client zu schützen}
}
\glsadd{gls-Server-Side-Token-Pattern}

\newglossaryentry{gls-Exponentielles-Backoff}{
  type=fachwort,
  name=Exponentielles Backoff,
  description={Strategie zur automatischen Wiederholung fehlgeschlagener Vorgänge mit exponentiell steigenden Wartezeiten}
}
\glsadd{gls-Exponentielles-Backoff}

\newglossaryentry{gls-Python-Package-Index}{
  type=fachwort,
  name=Python Package Index,
  description={Zentrales Repository für Python-Softwarepakete (auch PyPI genannt)}
}
\glsadd{gls-Python-Package-Index}

\newglossaryentry{gls-Zustandsmaschine}{
  type=fachwort,
  name=Zustandsmaschine,
  description={Modell zur Beschreibung des Verhaltens eines Systems durch definierte Zustände und Übergänge}
}
\glsadd{gls-Zustandsmaschine}

\newglossaryentry{gls-WebSocket}{
  type=fachwort,
  name=WebSocket,
  description={Protokoll für bidirektionale, persistente Kommunikation zwischen Client und Server}
}
\glsadd{gls-WebSocket}

\newglossaryentry{gls-Silero-VAD}{
  type=fachwort,
  name=Silero VAD,
  description={Lokal ausführbares Voice Activity Detection Modell zur Erkennung von Sprachaktivität in Audio}
}
\glsadd{gls-Silero-VAD}

\newglossaryentry{gls-GPT-4o-mini}{
  type=fachwort,
  name=GPT-4o-mini,
  description={Kompaktes Large Language Model von OpenAI mit reduziertem Ressourcenbedarf}
}
\glsadd{gls-GPT-4o-mini}

\newglossaryentry{gls-Docker-Container}{
  type=fachwort,
  name=Docker-Container,
  description={Laufende Instanz eines Docker-Images, die eine isolierte Anwendungsumgebung bereitstellt}
}
\glsadd{gls-Docker-Container}

\newglossaryentry{gls-Port-Forwarding}{
  type=fachwort,
  name=Port-Forwarding,
  description={Netzwerktechnik zur Weiterleitung eingehender Verbindungen von einem Port auf einen anderen Port oder ein anderes Gerät, auch als Portweiterleitung bezeichnet}
}
\glsadd{gls-Port-Forwarding}

\newglossaryentry{gls-Portfreigabe}{
  type=fachwort,
  name=Portfreigabe,
  description={Konfiguration an Router oder Firewall zur Freigabe eines Ports für eingehende Verbindungen aus dem Internet an ein internes Gerät}
}
\glsadd{gls-Portfreigabe}

\newglossaryentry{gls-Firewall}{
  type=fachwort,
  name=Firewall,
  description={Sicherheitssystem zur Kontrolle des ein- und ausgehenden Netzwerkverkehrs basierend auf definierten Regeln}
}
\glsadd{gls-Firewall}

\newglossaryentry{gls-ulaw}{
  type=fachwort,
  name=ulaw,
  description={Audio-Codec für Sprachkompression in der Telefonie (µ-law PCM), standardisiert für nordamerikanische Telefonnetze}
}
\glsadd{gls-ulaw}

\newglossaryentry{gls-alaw}{
  type=fachwort,
  name=alaw,
  description={Audio-Codec für Sprachkompression in der Telefonie (A-law PCM), standardisiert für europäische und internationale Telefonnetze}
}
\glsadd{gls-alaw}

\newglossaryentry{gls-Opus-Codec}{
  type=fachwort,
  name=Opus,
  description={Hocheffizienter Audio-Codec für Sprach- und Musikübertragung mit niedriger Latenz, standardisiert in RFC 6716 und weit verbreitet in WebRTC}
}
\glsadd{gls-Opus-Codec}

\newglossaryentry{gls-G711}{
  type=fachwort,
  name=G.711,
  description={ITU-T Standard für Sprachcodec mit PCM-Kodierung, weit verbreitet in klassischer Telefonie mit 64 kbit/s Datenrate}
}
\glsadd{gls-G711}

\newglossaryentry{gls-Catch-all-Fallback}{
  type=fachwort,
  name=Catch-all-Fallback,
  description={Regelkonfiguration zur Behandlung aller nicht explizit definierten Fälle in einem Routing- oder Verarbeitungssystem}
}
\glsadd{gls-Catch-all-Fallback}

\newglossaryentry{gls-SIP-INVITE}{
  type=fachwort,
  name=SIP INVITE,
  description={SIP-Nachrichtentyp zum Initiieren einer Session oder Einladung zu einem Anruf}
}
\glsadd{gls-SIP-INVITE}

\newglossaryentry{gls-Dispatch-Rule}{
  type=fachwort,
  name=Dispatch Rule,
  description={Regelwerk zur Weiterleitung eingehender SIP-Anrufe an bestimmte Ziele oder Räume}
}
\glsadd{gls-Dispatch-Rule}

\newglossaryentry{gls-Cleanup-Funktion}{
  type=fachwort,
  name=Cleanup-Funktion,
  description={Programmiermuster oder Funktion zur Freigabe von Ressourcen und Bereinigung von Systemzuständen}
}
\glsadd{gls-Cleanup-Funktion}

\newglossaryentry{gls-ISO-8601}{
  type=fachwort,
  name=ISO 8601,
  description={Internationaler Standard für die Darstellung von Datum und Uhrzeit in einheitlichem Format}
}
\glsadd{gls-ISO-8601}

\newglossaryentry{gls-Turn}{
  type=fachwort,
  name=Turn,
  description={Einzelner Dialog-Schritt in einer Konversation, bestehend aus Nutzeräußerung und Systemantwort}
}
\glsadd{gls-Turn}

\newglossaryentry{gls-Word-Error-Rate}{
  type=fachwort,
  name=Word-Error-Rate,
  description={Metrik zur Bewertung der Genauigkeit von Spracherkennung, die den Anteil falsch erkannter Wörter misst}
}
\glsadd{gls-Word-Error-Rate}

\newglossaryentry{gls-Client-Initiated}{
  type=fachwort,
  name=Client-Initiated,
  description={Verbindungsabbruch oder Aktion, die vom Client ausgelöst wurde}
}
\glsadd{gls-Client-Initiated}

\newglossaryentry{gls-Paketverlust}{
  type=fachwort,
  name=Paketverlust,
  description={Verlust von Datenpaketen während der Netzwerkübertragung, der zu Qualitätsminderung führen kann}
}
\glsadd{gls-Paketverlust}

\newglossaryentry{gls-Jitter}{
  type=fachwort,
  name=Jitter,
  description={Variabilität in der Verzögerung von Netzwerkpaketen, die zu unregelmäßiger Wiedergabe führen kann}
}
\glsadd{gls-Jitter}

\newglossaryentry{gls-Echo}{
  type=fachwort,
  name=Echo,
  description={Akustisches Phänomen, bei dem der eigene Ton mit Verzögerung zurückgehört wird}
}
\glsadd{gls-Echo}

\newglossaryentry{gls-Payload}{
  type=fachwort,
  name=Payload,
  description={Nutzdaten einer Netzwerknachricht oder API-Anfrage, die die eigentlichen Informationen enthalten}
}
\glsadd{gls-Payload}

\newglossaryentry{gls-Snapshot}{
  type=fachwort,
  name=Snapshot,
  description={Momentaufnahme eines Systemzustands zu einem bestimmten Zeitpunkt, die für Backup, Wiederherstellung oder Versionierung verwendet wird}
}
\glsadd{gls-Snapshot}

\newglossaryentry{gls-Incident-Response-Plan}{
  type=fachwort,
  name=Incident-Response-Plan,
  description={Dokumentierter Prozess zur Reaktion auf Sicherheitsvorfälle und Systemausfälle}
}
\glsadd{gls-Incident-Response-Plan}

\newglossaryentry{gls-Auto-Scaling}{
  type=fachwort,
  name=Auto-Scaling,
  description={Automatische Anpassung von Ressourcen (z.B. Anzahl Container) basierend auf Systemlast}
}
\glsadd{gls-Auto-Scaling}

\newglossaryentry{gls-Load-Balancer}{
  type=fachwort,
  name=Load-Balancer,
  description={System zur Verteilung von Netzwerkverkehr auf mehrere Server zur Optimierung von Auslastung und Verfügbarkeit}
}
\glsadd{gls-Load-Balancer}

\newglossaryentry{gls-Edge-Computing}{
  type=fachwort,
  name=Edge-Computing,
  description={Bereitstellung von Rechen- und Speicherressourcen nahe am Ort der Datenentstehung zur Reduzierung von Latenzen}
}
\glsadd{gls-Edge-Computing}

\newglossaryentry{gls-Multimodal}{
  type=fachwort,
  name=Multimodal,
  description={Verarbeitung und Integration mehrerer Modalitäten (z.B. Text, Audio, Video) in einem System}
}
\glsadd{gls-Multimodal}

\newglossaryentry{gls-On-Device-Inferenz}{
  type=fachwort,
  name=On-Device-Inferenz,
  description={Ausführung von maschinellen Lernmodellen direkt auf einem Endgerät statt in der Cloud}
}
\glsadd{gls-On-Device-Inferenz}

\newglossaryentry{gls-Trace}{
  type=fachwort,
  name=Trace,
  description={Aufzeichnung des Ausführungspfads einer Anfrage durch ein verteiltes System zur Performance-Analyse}
}
\glsadd{gls-Trace}

\newglossaryentry{gls-Timeout}{
  type=fachwort,
  name=Timeout,
  description={Zeitüberschreitung, nach der eine Verbindung oder Operation abgebrochen wird}
}
\glsadd{gls-Timeout}

\newglossaryentry{gls-Latenz}{
  type=fachwort,
  name=Latenz,
  description={Zeitverzögerung zwischen Senden und Empfangen von Daten in einem System}
}
\glsadd{gls-Latenz}

\newglossaryentry{gls-Throughput}{
  type=fachwort,
  name=Throughput,
  description={Datenmenge, die pro Zeiteinheit durch ein System übertragen werden kann}
}
\glsadd{gls-Throughput}

\newglossaryentry{gls-Bitrate}{
  type=fachwort,
  name=Bitrate,
  description={Datenmenge pro Zeiteinheit bei der Übertragung digitaler Medien, gemessen in Bits pro Sekunde}
}
\glsadd{gls-Bitrate}

\newglossaryentry{gls-Abtastrate}{
  type=fachwort,
  name=Abtastrate,
  description={Anzahl der Abtastwerte pro Sekunde bei der Digitalisierung analoger Signale, gemessen in Hertz (z.B. 8 kHz, 16 kHz, 48 kHz)}
}
\glsadd{gls-Abtastrate}

\newglossaryentry{gls-Best-Effort}{
  type=fachwort,
  name=Best-Effort,
  description={Dienstmodell ohne Leistungsgarantien, bei dem das System sein Bestes versucht}
}
\glsadd{gls-Best-Effort}

\newglossaryentry{gls-Uptime}{
  type=fachwort,
  name=Uptime,
  description={Zeitdauer, während der ein System betriebsbereit und verfügbar ist}
}
\glsadd{gls-Uptime}

\newglossaryentry{gls-Failover}{
  type=fachwort,
  name=Failover,
  description={Automatisches Umschalten auf ein Backup-System bei Ausfall des primären Systems}
}
\glsadd{gls-Failover}

\newglossaryentry{gls-Single-Host-Deployment}{
  type=fachwort,
  name=Single-Host-Deployment,
  description={Bereitstellung eines Systems auf einem einzelnen physischen oder virtuellen Server}
}
\glsadd{gls-Single-Host-Deployment}

\newglossaryentry{gls-Anrufaufkommen}{
  type=fachwort,
  name=Anrufaufkommen,
  description={Anzahl der Anrufe in einem bestimmten Zeitraum}
}
\glsadd{gls-Anrufaufkommen}

\newglossaryentry{gls-Transkript}{
  type=fachwort,
  name=Transkript,
  description={Textuelle Darstellung gesprochener Sprache}
}
\glsadd{gls-Transkript}

\newglossaryentry{gls-Inferenz}{
  type=fachwort,
  name=Inferenz,
  description={Anwendung eines trainierten maschinellen Lernmodells zur Vorhersage oder Generierung von Ausgaben}
}
\glsadd{gls-Inferenz}

\newglossaryentry{gls-Quantisierung}{
  type=fachwort,
  name=Quantisierung,
  description={Reduzierung der Präzision von Modellparametern zur Verringerung von Speicher- und Rechenbedarf bei maschinellen Lernmodellen}
}
\glsadd{gls-Quantisierung}

\newglossaryentry{gls-CPaaS}{
  type=fachwort,
  name=CPaaS,
  first={Communications Platform as a Service (CPaaS)},
  text=CPaaS,
  description={Communications Platform as a Service, eine Cloud-basierte Plattform zur Bereitstellung von Kommunikationsdiensten}
}

\newglossaryentry{gls-mTLS}{
  type=fachwort,
  name=mTLS,
  first={Mutual TLS (mTLS)},
  text=mTLS,
  description={Mutual TLS, eine Erweiterung von TLS, bei der beide Seiten (Client und Server) sich gegenseitig mit Zertifikaten authentifizieren}
}

\newglossaryentry{gls-QoS}{
  type=fachwort,
  name=QoS,
  first={Quality of Service (QoS)},
  text=QoS,
  description={Quality of Service, ein Konzept zur Verwaltung von Netzwerk-Performance-Parametern wie Latenz, Bandbreite und Paketverlust}
}

\newglossaryentry{gls-WebRTC}{
  type=fachwort,
  name=WebRTC,
  first={Web Real-Time Communication (WebRTC)},
  text=WebRTC,
  description={Web Real-Time Communication, ein Standard für Echtzeitkommunikation über Browser und Apps}
}

\newglossaryentry{gls-VoIP}{
  type=fachwort,
  name=VoIP,
  first={Voice over IP (VoIP)},
  text=VoIP,
  description={Voice over IP, Technologie zur Übertragung von Sprachkommunikation über IP-basierte Netzwerke, bei der Audiosignale digitalisiert, paketvermittelt übertragen und am Ziel wieder in Sprache umgewandelt werden}
}

\newglossaryentry{gls-WAV}{
  type=fachwort,
  name=WAV,
  first={Waveform Audio File Format (WAV)},
  text=WAV,
  description={Waveform Audio File Format, ein Audio-Dateiformat zur Speicherung von Audiosignalen in digitaler Form, das häufig für die Speicherung von Musik und Sprache verwendet wird}
}
